\documentclass[11pt, oneside]{article} 
\usepackage{geometry}
\geometry{letterpaper} 
\usepackage{graphicx}
	
\usepackage{amssymb}
\usepackage{amsmath}
\usepackage{parskip}
\usepackage{color}
\usepackage{hyperref}

\graphicspath{{/Users/telliott/Dropbox/Github-math/figures/}}
% \begin{center} \includegraphics [scale=0.4] {gauss3.png} \end{center}

\title{Tangent-secant theorem}
\date{}

\begin{document}
\maketitle
\Large

%[my-super-duper-separator]

We look at proofs of two theorems about secants and tangents, based on similar triangles.

\begin{center} \includegraphics [scale=0.25] {TS1.png} \end{center}

\[ PA \cdot PA' = PB \cdot PB' \]

It is always helpful to consider the form with the ratios:

\[ \frac{PA}{PB} = \frac{PB'}{PA'} \]

This strongly suggests we look at similar triangles.

\emph{Proof}.

We will show that $\angle PBA = \angle PA'B'$.  We might be tempted to apply the theorem which we called the "alternate segment theorem", but notice we do not have a tangent at $B$.

If we did, we would find that $\angle PBA$ is divided into two angles, one cutting off the arc $AB$ and the other with an equal vertical angle cutting off the arc $BB'$.  This combination is the arc cut off by $A'$.

As an alternative preliminary, we showed previously that for any quadrilateral whose four vertices all lie on one circle (a cyclic quadrilateral), the opposing vertices have supplementary angles.  Opposing vertices add to $180^{\circ}$ because their arc segments add up to one whole circle.

In the figure, $\angle ABB'$ is supplementary to both $\angle PBA$ and $\angle PA'B'$.  It follows that the angles marked with black dots are equal.

\begin{center} \includegraphics [scale=0.25] {TS1.png} \end{center}

Therefore $\triangle PBA$ is similar to $\triangle PA'B$.  The similarity is such that

\[ \frac{PA}{PB} = \frac{PB'}{PA'} \]

$\square$

Given a point $P$ outside the circle, the external part of any secant times the entire secant is a constant.
\[ PA \cdot PA' = PB \cdot PB' \]

One curious thing about this theorem is that these triangles are similar, and nested, but flipped.

We can do a bit more.

\subsection*{Tangent-secant theorem}

\label{sec:tangent_secant_theorem}

\begin{center} \includegraphics [scale=0.25] {TS2b.png} \end{center}

Let the points $B$ and $B'$ approach each other to become one point, re-labeled as $T$.  Then $PT$ will be a tangent of the circle.  Previously we had
\[ PA \cdot PA' = PB \cdot PB' \]

Now we modify it slightly:

\[ PA \cdot PA' = PT \cdot PT \]
\[ PA \cdot PA' = PT^2 \]

This is the tangent-secant theorem.

We must have two similar triangles, $\triangle PAT$ and $\triangle PTA'$.  The whole angle at vertex $T$ must correspond to $\angle PAT$.

Now run the logic backward and write the proof.

Note:  we've seen this sort of reverse logic a few times already.  The method has a name!  It was called the method of "analysis" by Pappus (320 A.D.), see Posamentier (Introduction).

To do this explicitly, show that since $\angle PTA$ includes the tangent, it cuts arc $AT$ just as $\angle PA'T$ does.

Instead, we take advantage of the tangent point to do something new.

\emph{Proof}.

\begin{center} \includegraphics [scale=0.25] {TS2.png} \end{center}

Draw the diameter $QOT \perp PT$ and also draw $QA$.  Since they correspond to equal arcs, the angles at $Q$ and $A'$ are equal.

By Thales' circle theorem, $\angle QAT$ is right.  So $\angle Q$ is complementary to $\angle ATQ$.

But the diameter is perpendicular to the tangent.  So $\angle ATQ$ is complementary to $\angle PTA$.

It follows that $\angle PTA = \angle PA'T$.

$\triangle PAT$ is similar to $\triangle PTA'$ by AAA, since they also share the angle at $P$.  This gives

\[ \frac{PT}{PA} = \frac{PA'}{PT} \]

which can be rearranged to the statement of the theorem:

\[ PA \cdot PA' = PT^2 \]

$\square$

One application of this theorem is to a determination of the size of the earth.

\begin{center} \includegraphics [scale=0.5] {al_biruni.png} \end{center} 

In the figure, the circle is the earth, of radius $R$, $h$ is the height of a convenient mountain, and $D$ is the distance to the horizon, which is tangent to the earth's radius.

Recall from the tangent-secant theorem 
\[ D^2 = h(2R + h) \]

We neglect $h^2$ compared to the other term so
\[ D^2 \approx 2Rh \]

About 1019 C.E., finding $h$ and $D$, Al-Biruni computed a value for $R$ equivalent to 3939 miles.

\end{document}